\section{Experiment}
\label{sec:simulation}

This section evaluates the proposed \tool framework under the probabilistic PMC setting. The evaluation is organized around three research questions:
\begin{itemize}[leftmargin=*, itemsep=1pt, parsep=0pt, topsep=1pt]
    \item \textbf{RQ1:} How effective is \tool compared with existing diagnosis methods?
    \item \textbf{RQ2:} How does \tool perform under different fault rates?
    \item \textbf{RQ3:} How do diagnostic features, posterior refinement, and GNN backbones affect \tool?
\end{itemize}

\subsection{Experimental Setup}

\noindent \textbf{Network topologies.}
We evaluate all methods on three representative interconnection topologies. The first topology is the hypercube $Q_n$, where each node is represented by an $n$-bit binary string and adjacent nodes differ in exactly one bit. The second topology is the $k$-ary $n$-cube $Q_k^n$, where each node is represented by an $n$-tuple over $\{0,1,\ldots,k-1\}$ and two nodes are adjacent if they differ by one modulo $k$ in exactly one dimension. The third topology is the augmented $k$-ary $n$-cube $AQ_{n,k}$, which extends $Q_k^n$ with additional augmented edges to provide higher connectivity. These three topologies cover different degrees, local connectivity patterns, and diagnostic evidence densities.

\noindent \textbf{Fault and syndrome generation.}
For each topology, faulty nodes are sampled according to a specified fault rate. Given the fault state of each node, directed test outcomes are generated according to the probabilistic PMC model. Each adjacent pair is tested in both directions, and each directed test is repeated for $T$ rounds unless otherwise specified. The proposed method uses the multi-round syndrome observations, while the rule-based baseline is evaluated using the last collected syndrome.

\noindent \textbf{Baselines.}
To the best of our knowledge, \tool is the first neural-network-based method for system-level probabilistic PMC diagnosis model. Therefore, we compare it with a representative traditional method, Clustered-PMC, the probabilistic clustered-fault diagnosis algorithm based on local factions and threshold $k=3$~\cite{sun2023probabilistic}. Since Clustered-PMC takes a single PMC syndrome as input, we apply it to the last collected syndrome in each test instance. Clustered-PMC is directly evaluated on the validation or test instances without training.

\noindent \textbf{Metrics.}
For binary diagnosis, we report accuracy, precision, recall, and F1-score. Since \tool produces posterior fault probabilities, we also report Brier score and expected calibration error (ECE) to evaluate probability quality. Lower Brier score and lower ECE indicate better calibrated posterior estimates.

\noindent \textbf{Implementation details.}
In RQ1 and RQ2, the GNN posterior estimator is instantiated with GraphSAGE. In RQ3, we use GraphSAGE for the feature and posterior-refinement studies, and separately compare GCN, GraphSAGE, and GAT in the backbone study. The default decision threshold for calibrated posterior diagnosis is $\theta_c=0.5$. The least-confident top-$q$ fraction in each preliminary class is selected for posterior refinement, with $q=10\%$ by default. The concrete values of $T$, fault rates, graph sizes, training epochs, hidden dimensions, and optimization settings are reported together with the corresponding experiments.

\noindent \textbf{Experimental environment.}
All experiments are conducted on a workstation equipped with an Intel Core i7-14700KF CPU and an NVIDIA GeForce RTX 5060 Ti GPU with 16 GB memory. The implementation uses Python 3.13, PyTorch 2.12.0 with CUDA 13.0 support, and NumPy 2.4.6. 

\subsection{RQ1: Overall Diagnostic Effectiveness}
\label{subsec:rq1}

RQ1 evaluates whether \tool can accurately diagnose faulty processors under the probabilistic PMC model. We compare \tool with Clustered-PMC on the three topologies described above. In the GPU-scale setting, we evaluate three larger instances from each topology family: $Q_6$, $Q_8$, $Q_{10}$, $Q_8^2$, $Q_6^3$, $Q_5^4$, $AQ_{2,8}$, $AQ_{3,6}$, and $AQ_{4,5}$. The fault rate is fixed to $30\%$, and each result is averaged over three independent runs.

For each topology instance and each run, we generate $10$ training samples and $5$ test samples with $T=2$ rounds of mutual testing. Unless otherwise stated, the probabilistic PMC parameters are sampled from $\gamma \sim U(0.7, 1.0)$ and $\beta \sim U(0.0, 0.3)$. The posterior estimator uses a two-layer GraphSAGE backbone with hidden dimension $64$, and is trained for $20$ epochs on GPU. The refinement candidate ratio is fixed to $q=10\%$. To improve training throughput for the GPU-scale setting, we use mini-batches of $16$ samples together with CUDA mixed-precision training. This low-data budget is aligned with the RQ3 ablation study so that the main comparisons and the component analysis are conducted under the same training constraint.

\begin{table}[h]
\centering
\caption{RQ1 diagnostic performance in the GPU-scale setting.}
\label{tab:rq1_effectiveness}
\resizebox{0.48\textwidth}{!}{%
\begin{tabular}{|c|c|c|c|c|c|c|}
\hline
\textbf{Topology} & \textbf{Method} & \textbf{Accuracy} & \textbf{Precision} & \textbf{Recall} & \textbf{F1} & \textbf{Brier} \\
\hline
$Q_6$ & Clustered-PMC & 0.8531 & 1.0000 & 0.5052 & 0.6403 & 0.1469 \\
$Q_6$ & \tool & 0.9937 & 1.0000 & 0.9789 & 0.9875 & 0.0238 \\
\hline
$Q_8$ & Clustered-PMC & 0.8114 & 1.0000 & 0.3732 & 0.4911 & 0.1886 \\
$Q_8$ & \tool & 0.9951 & 1.0000 & 0.9835 & 0.9906 & 0.0260 \\
\hline
$Q_{10}$ & Clustered-PMC & 0.7817 & 1.0000 & 0.2719 & 0.3740 & 0.2183 \\
$Q_{10}$ & \tool & 0.9920 & 1.0000 & 0.9733 & 0.9845 & 0.0298 \\
\hline
$Q_8^2$ & Clustered-PMC & 0.9031 & 0.9748 & 0.6947 & 0.7834 & 0.0969 \\
$Q_8^2$ & \tool & 0.9833 & 0.9868 & 0.9579 & 0.9706 & 0.0329 \\
\hline
$Q_6^3$ & Clustered-PMC & 0.8559 & 1.0000 & 0.5210 & 0.6663 & 0.1441 \\
$Q_6^3$ & \tool & 0.9997 & 0.9990 & 1.0000 & 0.9995 & 0.0240 \\
\hline
$Q_5^4$ & Clustered-PMC & 0.8254 & 1.0000 & 0.4195 & 0.5585 & 0.1746 \\
$Q_5^4$ & \tool & 0.9940 & 1.0000 & 0.9801 & 0.9896 & 0.0342 \\
\hline
$AQ_{2,8}$ & Clustered-PMC & 0.8396 & 0.9333 & 0.4597 & 0.5808 & 0.1604 \\
$AQ_{2,8}$ & \tool & 0.9948 & 0.9965 & 0.9859 & 0.9911 & 0.0323 \\
\hline
$AQ_{3,6}$ & Clustered-PMC & 0.8098 & 1.0000 & 0.3682 & 0.4615 & 0.1902 \\
$AQ_{3,6}$ & \tool & 1.0000 & 1.0000 & 1.0000 & 1.0000 & 0.0299 \\
\hline
$AQ_{4,5}$ & Clustered-PMC & 0.7473 & 1.0000 & 0.1599 & 0.2445 & 0.2527 \\
$AQ_{4,5}$ & \tool & 0.9908 & 1.0000 & 0.9695 & 0.9821 & 0.0308 \\
\hline
\end{tabular}%
}
\vspace{1mm}
\end{table}

Table~\ref{tab:rq1_effectiveness} shows that \tool consistently improves recall and F1-score over Clustered-PMC while maintaining high precision even under the same low-data budget used in the ablation study. Across all nine GPU-scale topologies, \tool achieves strong diagnosis quality, with F1 scores between $0.9706$ and $1.0000$ and Brier scores between $0.0238$ and $0.0342$. In contrast, Clustered-PMC remains much more conservative: its precision is usually close to one, but its recall drops substantially as the topology becomes larger or more densely augmented, leading to markedly lower F1 scores.

More specifically, on the hypercube family, the F1 score of Clustered-PMC decreases from $0.6403$ on $Q_6$ to $0.3740$ on $Q_{10}$, while \tool still maintains F1 values of $0.9875$, $0.9906$, and $0.9845$ on $Q_6$, $Q_8$, and $Q_{10}$, respectively. A similar pattern appears on the augmented $k$-ary family: Clustered-PMC reaches $0.5808$ on $AQ_{2,8}$, drops to $0.4615$ on $AQ_{3,6}$, and further drops to $0.2445$ on $AQ_{4,5}$, whereas \tool keeps F1 between $0.9821$ and $1.0000$. Even on the comparatively easier $k$-ary case $Q_8^2$, where Clustered-PMC attains its best F1 score of $0.7834$, \tool still improves the F1 score to $0.9706$. These results indicate that the posterior-estimation-and-refinement pipeline is not tied to one specific topology parameterization; instead, it remains highly effective across diverse $(k,n)$ combinations even when the available training data are very limited, while the rule-based baseline suffers mainly from low recall and therefore misses a nontrivial portion of faulty processors.

\subsection{RQ2: Performance Under Different Fault Rates}
\label{subsec:rq2_fault_rate}

RQ2 evaluates the robustness of \tool as the number of faulty processors increases. We vary the fault rate and compare \tool with Clustered-PMC under the same syndrome generation setting. This experiment is designed to test whether the posterior-based diagnosis remains effective when local neighborhoods contain more faulty nodes and the observed syndromes become more ambiguous.

In the GPU-scale setting, we use one representative larger instance from each topology family, namely $Q_{10}$, $Q_4^5$, and $AQ_{5,4}$. The fault rate is varied from $10\%$ to $50\%$ in steps of $5\%$, i.e., $\{10\%, 15\%, 20\%, 25\%, 30\%, 35\%, 40\%, 45\%, 50\%\}$. The remaining settings are the same as in RQ1: $T=2$, $\gamma \sim U(0.7, 1.0)$, $\beta \sim U(0.0, 0.3)$, $10$ training samples and $5$ test samples per topology instance, a two-layer GraphSAGE posterior estimator with hidden dimension $64$, $20$ training epochs, and posterior refinement ratio $q=10\%$. Each rate-wise result is averaged over the three topologies and three independent runs per topology.

\begin{table}[h]
\centering
\caption{Diagnostic performance under different fault rates.}
\label{tab:rq2_fault_rate}
\resizebox{0.48\textwidth}{!}{%
\begin{tabular}{|c|c|c|c|c|c|c|}
\hline
\textbf{Fault Rate} & \textbf{Method} & \textbf{Accuracy} & \textbf{Precision} & \textbf{Recall} & \textbf{F1} & \textbf{ECE} \\
\hline
10\% & Clustered-PMC & 0.9268 & 0.9356 & 0.2656 & 0.3631 & -- \\
10\% & \tool & 1.0000 & 1.0000 & 1.0000 & 1.0000 & 0.0000 \\
\hline
15\% & Clustered-PMC & 0.8855 & 0.9533 & 0.2384 & 0.3289 & -- \\
15\% & \tool & 1.0000 & 1.0000 & 1.0000 & 1.0000 & 0.0000 \\
\hline
20\% & Clustered-PMC & 0.8534 & 0.9822 & 0.2679 & 0.3686 & -- \\
20\% & \tool & 1.0000 & 1.0000 & 1.0000 & 1.0000 & 0.0000 \\
\hline
25\% & Clustered-PMC & 0.8073 & 0.9844 & 0.2291 & 0.3263 & -- \\
25\% & \tool & 1.0000 & 1.0000 & 1.0000 & 1.0000 & 0.0000 \\
\hline
30\% & Clustered-PMC & 0.7748 & 0.9867 & 0.2486 & 0.3496 & -- \\
30\% & \tool & 1.0000 & 1.0000 & 1.0000 & 1.0000 & 0.0000 \\
\hline
35\% & Clustered-PMC & 0.7356 & 0.9943 & 0.2438 & 0.3450 & -- \\
35\% & \tool & 1.0000 & 1.0000 & 1.0000 & 1.0000 & 0.0000 \\
\hline
40\% & Clustered-PMC & 0.6949 & 0.9977 & 0.2380 & 0.3398 & -- \\
40\% & \tool & 1.0000 & 1.0000 & 1.0000 & 1.0000 & 0.0000 \\
\hline
45\% & Clustered-PMC & 0.6618 & 0.9952 & 0.2489 & 0.3528 & -- \\
45\% & \tool & 0.9999 & 0.9999 & 0.9999 & 0.9999 & 0.0001 \\
\hline
50\% & Clustered-PMC & 0.6193 & 0.9959 & 0.2389 & 0.3436 & -- \\
50\% & \tool & 0.9998 & 0.9998 & 0.9998 & 0.9998 & 0.0002 \\
\hline
\end{tabular}%
}
\vspace{1mm}
\end{table}

Table~\ref{tab:rq2_fault_rate} shows that \tool remains highly robust as the fault rate increases even under the same low-data training budget. From $10\%$ to $40\%$, \tool maintains essentially perfect diagnosis, with accuracy and F1 both equal to $1.0000$ after rounding. Even at the highest tested fault rates of $45\%$ and $50\%$, the performance drop is negligible: the F1 score remains $0.9999$ and $0.9998$, respectively, while the ECE stays below $2.5 \times 10^{-4}$.

By contrast, Clustered-PMC degrades substantially as the fault rate increases. Its accuracy decreases from $0.9268$ at $10\%$ faults to $0.6193$ at $50\%$ faults. More importantly, its recall remains consistently low across all rates, fluctuating only between about $0.23$ and $0.27$, which leads to poor F1 scores in the range of $0.3263$--$0.3686$ for most settings. This behavior indicates that the rule-based baseline is overly conservative under the probabilistic PMC setting: it keeps precision high, but misses many faulty processors once the local syndromes become noisier and more ambiguous. In contrast, the posterior-based diagnosis in \tool degrades much more gracefully and preserves both discrimination and calibration quality over the full fault-rate range. Combined with the RQ1 results, this shows that the complete posterior-estimation-and-refinement pipeline remains highly effective not only when the fault rate changes, but also when the amount of available training data is restricted to the same low-data regime used in RQ3.

\subsection{RQ3: Ablation Study}
\label{subsec:rq3_ablation}

RQ3 studies three design questions in \tool. First, we evaluate the contribution of each node feature at the raw GNN-output level, before posterior refinement is applied. Second, we measure the contribution of local syndrome-consistency posterior refinement by comparing the raw and refined outputs under the same experimental setting. Third, we compare different GNN backbones to check whether the framework is tied to a specific graph encoder.

\noindent\textbf{Feature study.}
We first study the node feature design by removing one feature at a time while keeping the GNN backbone unchanged and evaluating the raw GNN outputs before posterior refinement:
\begin{itemize}[leftmargin=*, itemsep=1pt, parsep=0pt, topsep=1pt]
    \item \textbf{w/o match ratio}: removes the average match ratio from the node features.
    \item \textbf{w/o mismatch ratio}: removes the average mismatch ratio from the node features.
    \item \textbf{w/o dispersion}: removes mismatch dispersion from the node features.
    \item \textbf{full \tool}: uses all diagnostic node features under the default GraphSAGE posterior estimator.
\end{itemize}

In the updated GPU setting, this ablation is conducted on the smaller representative topology instances $Q_6$, $Q_8^2$, and $AQ_{2,8}$ with fault rate $30\%$, $T=2$, $\gamma \sim U(0.7, 1.0)$, and $\beta \sim U(0.0, 0.3)$. For each topology instance and each run, we generate $10$ training samples and $5$ test samples. The posterior estimator uses a two-layer GraphSAGE backbone with hidden dimension $64$. To further shorten the rerun budget while retaining a visible convergence gap, we train for $20$ epochs, evaluate every $5$ epochs, and use mini-batches of $16$ samples together with CUDA mixed-precision training. Each variant is averaged over the three topologies and three independent runs per topology. Since this ablation is intended to isolate the contribution of node features at the GNN level, Table~\ref{tab:rq3_feature_detail} reports the raw predictions before local posterior refinement; the refinement step is intentionally excluded from this first comparison.


\begin{table*}[t]
\centering
\caption{Complete topology-wise pre-refinement results for the feature ablation study.}
\label{tab:rq3_feature_detail}
\resizebox{\textwidth}{!}{%
\begin{tabular}{|c|c|c|c|c|c|c|c|c|c|}
\hline
\textbf{Topology} & \textbf{Variant} & \textbf{Accuracy} & \textbf{Precision} & \textbf{Recall} & \textbf{F1} & \textbf{ECE} & \textbf{Best-F1 Epoch} & \textbf{Epoch to F1 $\geq 0.99$} & \textbf{Mean F1 over Epochs} \\
\hline
$Q_6$ & w/o match ratio & 0.7125 & 0.2667 & 0.0316 & 0.0547 & 0.3209 & 7.33 & N/A & 0.0844 \\
$Q_6$ & w/o mismatch ratio & 0.7344 & 0.3333 & 0.1053 & 0.1462 & 0.2872 & 7.33 & N/A & 0.0903 \\
$Q_6$ & w/o dispersion & 0.8073 & 0.5961 & 0.3544 & 0.4030 & 0.2640 & 8.67 & N/A & 0.2273 \\
$Q_6$ & full \tool & 0.8333 & 0.7333 & 0.4386 & 0.5155 & 0.2834 & 20.00 & N/A & 0.1031 \\
\hline
$Q_8^2$ & w/o match ratio & 0.8177 & 0.5937 & 0.3930 & 0.4538 & 0.2945 & 1.00 & N/A & 0.2129 \\
$Q_8^2$ & w/o mismatch ratio & 0.7458 & 0.5333 & 0.1439 & 0.2148 & 0.2920 & 7.33 & N/A & 0.0946 \\
$Q_8^2$ & w/o dispersion & 0.9427 & 0.9796 & 0.8281 & 0.8820 & 0.2334 & 13.67 & N/A & 0.3423 \\
$Q_8^2$ & full \tool & 0.7885 & 0.3298 & 0.2912 & 0.3082 & 0.2747 & 7.33 & N/A & 0.0616 \\
\hline
$AQ_{2,8}$ & w/o match ratio & 0.8688 & 1.0000 & 0.5579 & 0.6612 & 0.2767 & 20.00 & N/A & 0.1322 \\
$AQ_{2,8}$ & w/o mismatch ratio & 0.7073 & 0.1333 & 0.0140 & 0.0248 & 0.3007 & 7.33 & N/A & 0.0050 \\
$AQ_{2,8}$ & w/o dispersion & 0.8917 & 0.9930 & 0.6421 & 0.7265 & 0.2411 & 20.00 & N/A & 0.1453 \\
$AQ_{2,8}$ & full \tool & 0.8521 & 0.6667 & 0.5018 & 0.5583 & 0.2592 & 13.67 & N/A & 0.1117 \\
\hline
\end{tabular}%
}
\vspace{1mm}
\end{table*}

Table~\ref{tab:rq3_feature_detail} reports the complete topology-wise pre-refinement results for the feature ablation study under the more extreme setting of only $10$ training samples and $5$ test samples. Without the masking effect of posterior refinement, the raw GNN outputs now exhibit severe under-training and much higher variance. Across all twelve topology-variant combinations, the final F1 values range from $0.0248$ to $0.8820$, while ECE ranges from $0.2334$ to $0.3209$. None of the feature variants reaches F1 $\geq 0.99$ within the $20$-epoch budget.

Under this $10$-sample budget, the feature study no longer supports a stable ranking among feature subsets. The raw GraphSAGE estimator is simply too under-trained for the pre-refinement feature comparison to remain robust. For example, on $Q_8^2$ the full feature set yields final F1 $0.3082$, whereas removing dispersion unexpectedly reaches $0.8820$; on $AQ_{2,8}$ the worst variant becomes ``w/o mismatch ratio'' with final F1 only $0.0248$. These reversals indicate that, once the number of training graphs is reduced to this level, optimization noise dominates the pre-refinement feature ablation. Therefore, this regime is still useful for stress-testing the full pipeline, but it is too data-starved to support a reliable feature-importance conclusion from raw GNN outputs alone.

\noindent\textbf{Posterior refinement study.}
We then isolate the effect of the local syndrome-consistency posterior refinement. Here, both variants use exactly the same trained GraphSAGE posterior estimator and the same node features as in the feature study; the only difference is whether we report the raw GNN outputs directly or apply the local posterior-refinement step to selected low-confidence predictions.

This study uses the same smaller GPU topology set as the feature ablation, namely $Q_6$, $Q_8^2$, and $AQ_{2,8}$, together with fault rate $30\%$, $10$ training samples, $5$ test samples, $20$ training epochs, evaluation every $5$ epochs, mini-batch size $16$, and CUDA mixed-precision training. The GraphSAGE posterior estimator and all feature settings are kept unchanged so that the only difference is whether the local posterior-refinement step is applied. Therefore, this second study should be read as a direct follow-up to the first one: it starts from the same raw GNN outputs and quantifies how much the refinement stage improves them.

\begin{table*}[t]
\centering
\caption{Complete topology-wise results for the posterior refinement study.}
\label{tab:rq3_refinement_detail}
\resizebox{\textwidth}{!}{%
\begin{tabular}{|c|c|c|c|c|c|c|c|c|c|}
\hline
\textbf{Topology} & \textbf{Variant} & \textbf{Accuracy} & \textbf{Precision} & \textbf{Recall} & \textbf{F1} & \textbf{ECE} & \textbf{Best-F1 Epoch} & \textbf{Epoch to F1 $\geq 0.99$} & \textbf{Mean F1 over Epochs} \\
\hline
$Q_6$ & w/o posterior refinement & 0.8854 & 0.8471 & 0.6351 & 0.6686 & 0.2605 & 20.00 & N/A & 0.3163 \\
$Q_6$ & full \tool & 0.9896 & 0.9967 & 0.9684 & 0.9801 & 0.1467 & 16.67 & 9.67 & 0.7049 \\
\hline
$Q_8^2$ & w/o posterior refinement & 0.8281 & 0.6596 & 0.4281 & 0.4930 & 0.2740 & 13.67 & N/A & 0.1291 \\
$Q_8^2$ & full \tool & 0.9458 & 0.9873 & 0.8281 & 0.8837 & 0.1513 & 16.67 & N/A & 0.6035 \\
\hline
$AQ_{2,8}$ & w/o posterior refinement & 0.8448 & 0.6590 & 0.4842 & 0.5405 & 0.2541 & 13.67 & N/A & 0.1081 \\
$AQ_{2,8}$ & full \tool & 0.9917 & 1.0000 & 0.9719 & 0.9848 & 0.1221 & 20.00 & 1.00 & 0.6187 \\
\hline
\end{tabular}%
}
\vspace{1mm}
\end{table*}

Table~\ref{tab:rq3_refinement_detail} shows that posterior refinement becomes indispensable under the more extreme $10$-train/$5$-test budget. The improvement is visible on all three topologies and is especially clear on $Q_8^2$, where refinement improves accuracy from $0.8281$ to $0.9458$, F1 from $0.4930$ to $0.8837$, and ECE from $0.2740$ to $0.1513$.

The convergence statistics show an even clearer effect. On $Q_6$, refinement raises the mean F1 over epochs from $0.3163$ to $0.7049$. On $AQ_{2,8}$, it raises the mean F1 over epochs from $0.1081$ to $0.6187$ while improving the final F1 from $0.5405$ to $0.9848$. This indicates that the refinement step does more than polish the final prediction: under a severely data-starved training budget, it rescues an otherwise unstable posterior estimator and restores much of the final diagnosis quality. Therefore, the refinement stage is not merely a marginal improvement in this regime, but the main reason the full pipeline remains usable when only a handful of training graphs are available.

\noindent\textbf{Backbone comparison.}
We further instantiate the GNN posterior estimator with different graph encoders while keeping the same diagnostic features and posterior refinement procedure. The reported results are averaged over the three representative topologies.

This comparison uses the same smaller GPU topology set $Q_6$, $Q_8^2$, and $AQ_{2,8}$, with fault rate $30\%$, $10$ training samples, $5$ test samples, $20$ training epochs, evaluation every $5$ epochs, mini-batch size $16$, and CUDA mixed-precision training. All backbones use the same node features and the same posterior-refinement procedure so that the only varying component is the graph encoder.

\begin{table*}[t]
\centering
\caption{Complete topology-wise results for the backbone comparison study.}
\label{tab:rq3_backbone_detail}
\resizebox{\textwidth}{!}{%
\begin{tabular}{|c|c|c|c|c|c|c|c|c|c|}
\hline
\textbf{Topology} & \textbf{Backbone} & \textbf{Accuracy} & \textbf{Precision} & \textbf{Recall} & \textbf{F1} & \textbf{ECE} & \textbf{Best-F1 Epoch} & \textbf{Epoch to F1 $\geq 0.99$} & \textbf{Mean F1 over Epochs} \\
\hline
$Q_6$ & GCN & 0.8250 & 0.9810 & 0.4175 & 0.5770 & 0.2439 & 10.00 & 1.00 & 0.4902 \\
$Q_6$ & GraphSAGE & 0.9177 & 0.9905 & 0.7263 & 0.8075 & 0.1642 & 16.67 & 6.00 & 0.5904 \\
$Q_6$ & GAT & 0.8490 & 0.9844 & 0.4982 & 0.6371 & 0.1643 & 8.33 & 4.67 & 0.4722 \\
\hline
$Q_8^2$ & GCN & 0.8469 & 0.9506 & 0.5158 & 0.6361 & 0.2258 & 13.33 & N/A & 0.5058 \\
$Q_8^2$ & GraphSAGE & 0.8708 & 0.9679 & 0.5895 & 0.6931 & 0.2206 & 13.33 & N/A & 0.5648 \\
$Q_8^2$ & GAT & 0.8354 & 0.9775 & 0.4561 & 0.6071 & 0.1133 & 7.00 & N/A & 0.5002 \\
\hline
$AQ_{2,8}$ & GCN & 0.8250 & 0.9905 & 0.4140 & 0.5795 & 0.2315 & 10.00 & N/A & 0.4701 \\
$AQ_{2,8}$ & GraphSAGE & 0.9958 & 1.0000 & 0.9860 & 0.9927 & 0.1291 & 20.00 & 13.00 & 0.7033 \\
$AQ_{2,8}$ & GAT & 0.8083 & 0.9810 & 0.3614 & 0.5282 & 0.2611 & 10.00 & N/A & 0.4356 \\
\hline
\end{tabular}%
}
\vspace{1mm}
\end{table*}

Table~\ref{tab:rq3_backbone_detail} shows that all three graph encoders remain trainable under the $10$-sample budget, but their absolute performance drops sharply. On $Q_6$, GraphSAGE still attains the best final F1 ($0.8075$), ahead of GAT ($0.6371$) and GCN ($0.5770$). On $Q_8^2$, all three backbones degrade substantially, with GraphSAGE still leading at F1 $0.6931$ while GCN and GAT achieve $0.6361$ and $0.6071$, respectively. On $AQ_{2,8}$, GraphSAGE remains clearly more robust and still reaches final F1 $0.9927$, whereas GCN and GAT fall to $0.5795$ and $0.5282$. The calibration gaps also become large for all three encoders, indicating that this regime stresses both discrimination and confidence estimation.

Taken together, Tables~\ref{tab:rq3_feature_detail}, \ref{tab:rq3_refinement_detail}, and~\ref{tab:rq3_backbone_detail} support three observations on the smaller topology set under the $10$-train/$5$-test stress setting. First, the raw feature ablation becomes too unstable to support a reliable feature-importance conclusion, because the posterior estimator itself is severely under-trained. Second, posterior refinement becomes the dominant contributor to practical diagnosis quality, recovering large portions of the lost F1 and substantially reducing calibration error on all three topologies. Third, among the tested GNN encoders, GraphSAGE remains the most robust backbone, but the entire comparison also shows that $10$ training graphs place the learning problem close to the edge of failure on the harder topologies. Therefore, this regime is best interpreted as a stress test that highlights the necessity of posterior refinement and the relative robustness of GraphSAGE, rather than as the main setting for drawing fine-grained feature-ablation conclusions.
